%-------------------------
% Stepanenko Nazar — Embedded / Miltech Software Engineer CV
% Based on Jake's Resume template (https://github.com/jakegut/resume)
% License: MIT
%-------------------------

\documentclass[letterpaper,11pt]{article}

\usepackage{latexsym}
\usepackage[empty]{fullpage}
\usepackage{titlesec}
\usepackage{marvosym}
\usepackage[usenames,dvipsnames]{color}
\usepackage{verbatim}
\usepackage{enumitem}
\usepackage[hidelinks]{hyperref}
\usepackage{fancyhdr}
\usepackage[english]{babel}
\usepackage{tabularx}
\usepackage{fontawesome5}
\usepackage{multicol}
\usepackage{xcolor}
\setlength{\multicolsep}{-3.0pt}
\setlength{\columnsep}{-1pt}
\input{glyphtounicode}

\pagestyle{fancy}
\fancyhf{}
\fancyfoot{}
\renewcommand{\headrulewidth}{0pt}
\renewcommand{\footrulewidth}{0pt}

\addtolength{\oddsidemargin}{-0.6in}
\addtolength{\evensidemargin}{-0.5in}
\addtolength{\textwidth}{1.19in}
\addtolength{\topmargin}{-.7in}
\addtolength{\textheight}{1.4in}

\urlstyle{same}

\raggedbottom
\raggedright
\setlength{\tabcolsep}{0in}

% Accent color
\definecolor{accent}{HTML}{1F6FB2}

% Sections formatting
\titleformat{\section}{
  \vspace{-4pt}\scshape\raggedright\large\bfseries\color{accent}
}{}{0em}{}[\color{accent}\titlerule \vspace{-5pt}]

% Ensure that generate pdf is machine readable/ATS parsable
\pdfgentounicode=1

%-------------------------
% Custom commands
\newcommand{\resumeItem}[1]{
  \item\small{
    {#1 \vspace{-2pt}}
  }
}

\newcommand{\resumeSubheading}[4]{
  \vspace{-2pt}\item
    \begin{tabular*}{0.97\textwidth}[t]{l@{\extracolsep{\fill}}r}
      \textbf{#1} & #2 \\
      \textit{\small#3} & \textit{\small #4} \\
    \end{tabular*}\vspace{-7pt}
}

\newcommand{\resumeSubSubheading}[2]{
    \item
    \begin{tabular*}{0.97\textwidth}{l@{\extracolsep{\fill}}r}
      \textit{\small#1} & \textit{\small #2} \\
    \end{tabular*}\vspace{-7pt}
}

\newcommand{\resumeProjectHeading}[2]{
    \item
    \begin{tabular*}{0.97\textwidth}{l@{\extracolsep{\fill}}r}
      \small#1 & #2 \\
    \end{tabular*}\vspace{-7pt}
}

\newcommand{\resumeSubItem}[1]{\resumeItem{#1}\vspace{-4pt}}

\renewcommand\labelitemii{$\vcenter{\hbox{\tiny$\bullet$}}$}

\newcommand{\resumeSubHeadingListStart}{\begin{itemize}[leftmargin=0.15in, label={}]}
\newcommand{\resumeSubHeadingListEnd}{\end{itemize}}
\newcommand{\resumeItemListStart}{\begin{itemize}}
\newcommand{\resumeItemListEnd}{\end{itemize}\vspace{-5pt}}

%-------------------------------------------
%%%%%%  RESUME STARTS HERE  %%%%%%%%%%%%%%%%
%-------------------------------------------

\begin{document}

%----------HEADING----------
\begin{center}
    {\Huge \scshape Stepanenko Nazar} \\ \vspace{4pt}
    \small {\color{accent}\textbf{Embedded / Miltech Software Engineer (Junior)}} \\ \vspace{4pt}
    \faMapMarker* \, Kyiv, Ukraine \quad
    \faEnvelope \, \href{mailto:gangsterstepa@gmail.com}{gangsterstepa@gmail.com} \quad
    \faGithub \, \href{https://github.com/d3Par1}{github.com/d3Par1} \quad
    \faLinkedin \, \href{https://www.linkedin.com/in/yourhandle}{LinkedIn} \quad
    \faTelegram \, @yourhandle
\end{center}

%-----------SUMMARY-----------
\section{Summary}
\small
Computer Science student (2nd year, Igor Sikorsky KPI) focused on embedded systems, Linux, and unmanned platform software for defense applications. Strong foundation in C/C++ and systems programming, hands-on practice with Raspberry Pi Pico W and ESP32. Building experience toward STM32 bare-metal, FreeRTOS, and drone software stacks (PX4 / MAVLink) to contribute to Ukrainian miltech and autonomous-systems engineering.
\vspace{-6pt}

%-----------SKILLS-----------
\section{Technical Skills}
\begin{itemize}[leftmargin=0.15in, label={}]
    \small{\item{
     \textbf{Programming Languages}{: C / C++, Python, C\#, JavaScript / Node.js, Lua, SQL, HTML/CSS} \\
     \textbf{Embedded \& Systems}{: Raspberry Pi Pico W (hands-on), ESP32 (Wi-Fi/BLE), Linux (CLI, Bash), GPIO / UART / I\textsuperscript{2}C / SPI basics, RTOS concepts (learning), Real-time data \& sensor integration} \\
     \textbf{Tools \& Workflow}{: Git / GitHub, VS Code, CLion, CMake / Makefiles, PlatformIO, Arduino IDE, Docker (basic), Wireshark (basic), Blender} \\
     \textbf{Domain Interest}{: UAV / drone software, autonomous systems, sensor fusion, defense electronics, secure firmware}
    }}
\end{itemize}
\vspace{-16pt}

%-----------LANGUAGES-----------
\section{Spoken Languages}
\begin{itemize}[leftmargin=0.15in, label={}]
    \small{\item
        Ukrainian (Native) \quad\textbullet\quad
        English (Advanced \textbf{C1}) \quad\textbullet\quad
        French (Intermediate \textbf{B1--B2}) \quad\textbullet\quad
        German (Beginner \textbf{A1--A2})
    }
\end{itemize}
\vspace{-16pt}

%-----------PROJECTS-----------
\section{Projects}
\resumeSubHeadingListStart

  \resumeProjectHeading
    {\textbf{Microcontroller Hands-On Practice} $|$ \emph{C/C++, MicroPython, Arduino IDE, PlatformIO}}{2025--Present}
    \resumeItemListStart
      \resumeItem{Self-driven exploratory projects on Raspberry Pi Pico W and ESP32: GPIO control, sensor reading, Wi-Fi connectivity, basic peripherals (LEDs, buttons, OLED, sensors).}
      \resumeItem{Embedded development workflow practice: toolchain setup, flashing, serial debugging, datasheet reading.}
      \resumeItem{Currently extending toward STM32 bare-metal (HAL/CMSIS) and FreeRTOS.}
    \resumeItemListEnd

  \resumeProjectHeading
    {\textbf{Personal Finance Manager} $|$ \emph{Node.js, Express, SQLite, Bootstrap, Chart.js}}{2026}
    \resumeItemListStart
      \resumeItem{Built full-stack web app with auth (bcrypt + sessions), multi-account architecture, REST API for transactions.}
      \resumeItem{Designed interactive analytics dashboard with Chart.js, server-side aggregation via SQL \texttt{GROUP BY}.}
      \resumeItem{Team project (2 engineers); owned authentication, accounts, dashboard, and API layers end-to-end.}
      \resumeItem{Repo: \href{https://github.com/d3Par1/Coursework_Web}{github.com/d3Par1/Coursework\_Web}}
    \resumeItemListEnd

  \resumeProjectHeading
    {\textbf{Telegram Bots} $|$ \emph{Python, aiogram / python-telegram-bot, SQLite}}{2025--Present}
    \resumeItemListStart
      \resumeItem{Built several Telegram bots exploring command handlers, inline keyboards, FSM, and webhook deployment.}
      \resumeItem{Practiced async Python (\texttt{asyncio}), HTTP integration with external APIs, persistent state via SQLite.}
      \resumeItem{End-to-end ownership of small backend products from requirements to deployment.}
    \resumeItemListEnd

  \resumeProjectHeading
    {\textbf{University Coursework — System Software \& OOP} $|$ \emph{C, C++, C\#}}{2025--2026}
    \resumeItemListStart
      \resumeItem{Implemented low-level system programs in C: memory management, process scheduling, IPC primitives.}
      \resumeItem{Designed object-oriented systems in C\# covering SOLID principles, design patterns, software architecture.}
      \resumeItem{Repo: \href{https://github.com/d3Par1/University-Repository}{github.com/d3Par1/University-Repository}}
    \resumeItemListEnd

  \resumeProjectHeading
    {\textbf{Roblox Grinding Simulator} $|$ \emph{Lua, Vide reactive UI, Blender}}{2025--Present}
    \resumeItemListStart
      \resumeItem{Solo-developed gameplay loop, inventory, shop systems demonstrating real-time state management.}
      \resumeItem{Designed adaptive UI with cross-device scaling (mobile/desktop), reactive UI architecture.}
      \resumeItem{Created custom 3D assets and preview scenes in Blender — full ownership of code, UI, visual content.}
    \resumeItemListEnd

\resumeSubHeadingListEnd
\vspace{-12pt}

%-----------LEARNING ROADMAP-----------
\section{In Progress (2026 Embedded Roadmap)}
\begin{itemize}[leftmargin=0.15in, label={\small\textbullet}]
    \small{
      \item STM32 bare-metal firmware (HAL/CMSIS): timers, interrupts, UART telemetry — telemetry transmitter prototype
      \item FreeRTOS task scheduling: multi-task firmware with mutexes, queues, shared peripherals
      \item MAVLink protocol: message parsing in C; integration with PX4 / ArduPilot SITL
      \item Drone simulator (PX4 + Gazebo): autonomous mission scripting in Python / C++
      \item ROS2 fundamentals: publisher / subscriber nodes, sensor data pipelines
    }
\end{itemize}
\vspace{-12pt}

%-----------EDUCATION-----------
\section{Education}
\resumeSubHeadingListStart
  \resumeSubheading
    {Igor Sikorsky Kyiv Polytechnic Institute (KPI)}{Sep. 2024 -- Jun. 2028 (expected)}
    {B.Sc. Software Engineering \, $|$ \, Group ТВ-43}{Kyiv, Ukraine}
    \resumeItemListStart
      \resumeItem{\textbf{Relevant coursework:} System Software Architecture, Object-Oriented Analysis \& Design, Web Programming, Data Structures \& Algorithms, Software Engineering Components, Operating Systems.}
    \resumeItemListEnd
\resumeSubHeadingListEnd
\vspace{-12pt}

%-----------TRAININGS-----------
\section{Trainings \& Self-Directed Learning}
\begin{itemize}[leftmargin=0.15in, label={\small\textbullet}]
    \small{
      \item \textbf{Linux Basics} — KodeKloud (online, 2024). Shell, file system, processes, networking basics.
      \item \textbf{Web Programming Foundations} — KPI structured course (2025--2026). HTTP, REST, SSR, RDBMS.
      \item \textbf{Telegram Bot API \& async Python} — self-paced, multiple side projects (2025--present).
      \item \textbf{Blender 3D modeling} — self-paced via official tutorials and community resources.
      \item \textbf{Microcontroller hands-on} — self-paced experimentation with Raspberry Pi Pico W and ESP32.
    }
\end{itemize}
\vspace{-12pt}

%-----------ADDITIONAL-----------
\section{Additional}
\begin{itemize}[leftmargin=0.15in, label={\small\textbullet}]
    \small{
      \item \textbf{Defense-tech motivation:} Ukrainian engineer committed to contributing technical capability to defense and dual-use systems.
      \item \textbf{Self-learner:} Comfortable picking up new languages, frameworks, hardware platforms from official documentation.
      \item \textbf{Open to:} Internship $\cdot$ Working Student $\cdot$ Junior roles (Kyiv / Remote / Hybrid).
    }
\end{itemize}

\end{document}
